\ifdefined\gkver\else\def\gkver{student}\fi
\documentclass[\gkver]{gaokaozhenti}
\begin{document}

\chapter{2014年福建理综}
%% paperType: 理综物理部分

\begin{enumerate}
\item
%% number: 13
%% paperName: 2014年福建理综第 13 题 6 分
%% typeId: 1
%% type: 单选题
%% chapter: 光学
%% point: 光的折射
%% method: 无
%% score: 6
%% degree: 650
%% duplicateId: 0
%% body:
如图，一束光由空气射向半圆柱体玻璃砖，$O$ 点为该玻璃砖截面的圆心，下图能正确描述其光路图的是\xzanswer{A}

\fourchoices[answer=A, ispicture=true, h=2.20cm]
{figs/13a.jpg}
{figs/13b.jpg}
{figs/13c.jpg}
{figs/13d.jpg}
%% answer: A
%% memo:
\memoanswer{
当光由光疏介质射入光密介质，在界面同时发生反射和折射，且入射角大于折射角，B、D 项错误；当光由光密介质射入光疏介质时，如果入射角大于临界角就会发生全反射，如 A 选项所示，A 项正确；如果入射角小于临界角，同时发生反射和折射，且入射角小于折射角，C 项错误。
}

\item
%% number: 14
%% paperName: 2014年福建理综第 14 题 6 分
%% typeId: 1
%% type: 单选题
%% chapter: 曲线运动
%% point: 人造地球卫星
%% method: 无
%% score: 6
%% degree: 650
%% duplicateId: 0
%% body:
若有一颗“宜居”行星，其质量为地球的 $p$ 倍，半径为地球的 $q$ 倍，则该行星卫星的环绕速度是地球卫星环绕速度的\xzanswer{C}

\fourchoices[answer=C]
{$\sqrt {pq} $ 倍}
{$\sqrt {\frac{q}{p}} $ 倍}
{$\sqrt {\frac{p}{q}} $ 倍}
{$\sqrt {p{q^3}} $ 倍}
%% answer: C
%% memo:
\memoanswer{
第一宇宙速度又叫环绕速度，即绕星球表面飞行的卫星的速度，根据万有引力提供向心力 $\frac{GMm}{R^{2}} = \frac{mv^{2}}{R}$，解得 $v = \sqrt{\frac{GM}{R}}$，所以 $\frac{v}{v_{\text{地}}} = \sqrt{\frac{MR_{\text{地}}}{M_{\text{地}}R}} = \sqrt{\frac{p}{q}}$，C 项正确。
}

\item
%% number: 15
%% paperName: 2014年福建理综第 15 题 6 分
%% typeId: 1
%% type: 单选题
%% chapter: 直线运动
%% point: 运动图象
%% method: 无
%% score: 6
%% degree: 650
%% duplicateId: 0
%% body:
如右图，滑块以初速度 $v_{0}$ 沿表面粗糙且足够长的固定斜面，从顶端下滑，直至速度为零。对于该运动过程，若用 $h$、$s$、$v$、$a$ 分别表示滑块的下降高度、位移、速度和加速度的大小，$t$ 表示时间，则下列图像最能正确描述这一运动规律的是\xzanswer{B}

\onepicture[w=3.60cm, align=right]{figs/15.png}

\fourchoices[answer=B, ispicture=true, h=2.20cm]
{figs/15a.png}
{figs/15b.png}
{figs/15c.png}
{figs/15d.png}
%% answer: B
%% memo:
\memoanswer{
滑块沿斜面做匀减速直线运动，速度随时间均匀减小，加速度恒定，C、D 项错误；位移时间图像中斜率代表速度，A 项中斜率表示竖直方向的分速度在增加，A 项错误，B 项斜率表示合速度在减小，B 项正确。
}

\item
%% number: 16
%% paperName: 2014年福建理综第 16 题 6 分
%% typeId: 1
%% type: 单选题
%% chapter: 交流电
%% point: 变压器
%% method: 无
%% score: 6
%% degree: 450
%% duplicateId: 0
%% body:
图为模拟远距离输电实验电路图，两理想变压器的匝数 $n_{1} = n_{4} < n_{2} = n_{3}$，四根模拟输电线的电阻 $R_{1}$、$R_{2}$、$R_{3}$、$R_{4}$ 的阻值均为 $R$，$A_{1}$、$A_{2}$ 为相同的理想交流电流表，$L_{1}$、$L_{2}$ 为相同的小灯泡，灯丝电阻 $R_{L} > 2R$，忽略灯丝电阻随温度的变化，当 A、B 端接入低压交流电源时\xzanswer{D}

\onepicture[w=7.50cm]{figs/16.png}

\fourchoices[answer=D]
{$A_{1}$、$A_{2}$ 两表的示数相同}
{$L_{1}$、$L_{2}$ 两灯泡的亮度相同}
{$R_{1}$ 消耗的功率大于 $R_{3}$ 消耗的功率}
{$R_{2}$ 两端的电压小于 $R_{4}$ 两端的电压}
%% answer: D
%% memo:
\memoanswer{
由于 $n_{1} < n_{2}$，所以左边的理想变压器为升压变压器，低压交流电源经过理想升压变压器后，由于功率不变，所以电流变小，则 $A_{1}$ 表的示数小于 $A_{2}$ 表的示数，选项 A 错误；由 $\frac{U_{1}}{U_{2}} = \frac{n_{1}}{n_{2}}$，$U_{3} = U_{2} - I_{A1} \cdot 2R$ 及 $\frac{U_{3}}{U_{4}} = \frac{n_{3}}{n_{4}}$ 可知，$L_{1}$ 两端的电压 $U_{L1} = U_{4} = U_{1} - \frac{n_{4}}{n_{3}} \cdot I_{A1} \cdot 2R$，又 $L_{2}$ 两端的电压 $U_{L2} = U_{1} - I_{A2} \cdot 2R$，因 $\frac{n_{4}}{n_{3}} I_{A1} < I_{A2}$，故 $U_{L1} > U_{L2}$，由 $P = \frac{U^2}{R}$ 可知，$L_{1}$ 的功率大于 $L_{2}$ 的功率，所以 $L_{1}$ 的亮度大于 $L_{2}$，选项 B 错误；由于 $A_{1}$ 表的示数小于 $A_{2}$ 表的示数，由 $P = I^2 R$ 可知，$R_{1}$ 消耗的功率小于 $R_{3}$ 消耗的功率，选项 C 错误；由于 $A_{1}$ 表的示数小于 $A_{2}$ 表的示数，由 $U = IR$ 可知，$R_{2}$ 两端的电压小于 $R_{4}$ 两端的电压，选项 D 正确。
}

\item
%% number: 17
%% paperName: 2014年福建理综第 17 题 6 分
%% typeId: 1
%% type: 单选题
%% chapter: 机械波
%% point: 波的图象
%% method: 无
%% score: 6
%% degree: 650
%% duplicateId: 0
%% body:
在均匀介质中，一列沿 $x$ 轴正向传播的横波，其波源 $O$ 在第一个周期内的振动图像如右图所示，则该波在第一个周期末的波形图是\xzanswer{D}

\onepicture[w=3.60cm, align=right]{\includesvg{figs/17}}

\fourchoices[answer=D, ispicture=true, h=1.60cm]
{figs/17a.png}
{figs/17b.png}
{figs/17c.png}
{figs/17d.png}
%% answer: D
%% memo:
\memoanswer{
根据振动图像可知波源起振方向向下，则每一个点起振时方向均向下，当经历一个周期时，波源仍处于平衡位置且向下振动。由波动图象可知，沿着波的传播方向看：“上坡下，下坡上”，所以第一个周期末时波源处于上坡，波源完全重复上一个周期的运动，D 项正确。
}

\item
%% number: 18
%% paperName: 2014年福建理综第 18 题 6 分
%% typeId: 1
%% type: 单选题
%% chapter: 运动和力
%% point: 变加速直线运动
%% method: 无
%% score: 6
%% degree: 450
%% duplicateId: 0
%% body:
如图，两根相同的轻质弹簧，沿足够长的光滑斜面放置，下端固定在斜面底部挡板上，斜面固定不动。质量不同、形状相同的两物块分别置于两弹簧上端。现用外力作用在物块上，使两弹簧具有相同的压缩量，若撤去外力后，两物块由静止沿斜面向上弹出并离开弹簧，则从撤去外力到物块速度第一次减为零的过程，两物块\xzanswer{C}

\onepicture[w=4.20cm, align=right]{figs/18.jpg}

\fourchoices[answer=C]
{最大速度相同}
{最大加速度相同}
{上升的最大高度不同}
{重力势能的变化量不同}
%% answer: C
%% memo:
\memoanswer{
撤去外力后，两物块受到重力、弹力和支持力作用，在斜面上做往复运动，当弹簧弹力等于重力沿斜面向下的分力时，速度最大，根据牛顿第二定律和动能定理，有：$mg\sin\alpha=k\Delta x$，$W_{\text{合}}=\frac{1}{2}mv_{\text{max}}^{2}$，质量较小的物块平衡位置较高，合力做功较大，物块的最大速度较大，A 项错误；初位置时物块具有的加速度最大，由于弹簧的压缩量一定，所以最大弹力相同，根据牛顿第二定律有：$k\Delta x-mg\sin\alpha=ma$ 得 $a = \frac{k\Delta x}{m} - g\sin\alpha$，由于质量不同，所以最大加速度也不同，B 项错误；由于弹簧相同，两个物块的质量不同，所以速度最大时的位置不同，C 项正确；速度为零时，物块的动能也为零，减少的弹性势能全部转化为物块的重力势能，D 项错误。
}

\item
%% number: 19
%% paperName: 2014年福建理综第 19 题 12 分
%% typeId: 4
%% type: 实验题
%% chapter: 电路
%% point: 测电动势与内阻
%% method: 无
%% score: 12
%% degree: 450
%% duplicateId: 0
%% body:
\begin{enumerate}
	\item
	（6 分）某同学测定一金属杆的长度和直径，示数如图甲、乙所示，则该金属杆的长度和直径分别为\tkanswer[1.6cm]{6.010}$\Ucm$ 和\tkanswer[1.4cm]{4.20}$\Umm$。

	\twopicture[numstyle=chinese, w=5.00cm]{figs/19a.jpg}{figs/19b.jpg}
	\item
	（12 分）某研究性学习小组利用伏安法测定某一电池组的电动势和内阻，实验原理如图甲所示，其中，虚线框内为用灵敏电流计 G 改装的电流表 A，V 为标准电压表，$E$ 为待测电池组，S 为开关，$R$ 为滑动变阻器，$R_{0}$ 是标称值为 $4.0\UO$ 的定值电阻。

	\twopicture[numstyle=chinese, w=4.50cm]{figs/19c.jpg}{figs/19d.jpg}

	\begin{steps}
		\item
		已知灵敏电流计 G 的满偏电流 $I_{g} = 100\UuA$、内阻 $r_{g} = 2.0\UkO$，若要改装后的电流表满偏电流为 $200\UmA$，应并联一只\tkanswer[1.4cm]{1.0}$\UO$（保留一位小数）的定值电阻 $R_{1}$；
		\item
		根据图甲，用笔画线代替导线将图乙连接成完整电路；
		\item
		某次试验的数据如下表所示：

		\begin{tblr}{|c|c|c|c|c|c|c|c|c|}
		\hline
		测量次数 & 1 & 2 & 3 & 4 & 5 & 6 & 7 & 8 \\
		\hline
		电压表 V 读数 $U/\UV$ & 5.26 & 5.16 & 5.04 & 4.94 & 4.83 & 4.71 & 4.59 & 4.46 \\
		\hline
		改装表 A 读数 $I/\UmA$ & 20 & 40 & 60 & 80 & 100 & 120 & 140 & 160 \\
		\hline
		\end{tblr}

		该小组借鉴“研究匀变速直线运动”试验中计算加速度的方法（逐差法），计算出电池组的内阻 $r =\tkanswer[1.2cm]{1.66}\UO$（保留两位小数）；为减小偶然误差，逐差法在数据处理方面体现出的主要优点是\tkanswer[1.6cm]{充分利用已测得数据}。
		\item
		该小组在前面实验的基础上，为探究图甲电路中各元器件的实际阻值对测量结果的影响，用一已知电动势和内阻的标准电池组通过上述方法多次测量后发现：电动势的测量值与已知值几乎相同，但内阻的测量值总是偏大。若测量过程无误，则内阻测量值总是偏大的原因是\tkanswer[1.2cm]{CD}。（填选项前的字母）

		（A）电压表内阻的影响

		（B）滑动变阻器的最大阻值偏小

		（C）$R_{1}$ 的实际阻值比计算值偏小

		（D）$R_{0}$ 的实际阻值比标称值偏大
	\end{steps}
\end{enumerate}
%% answer: 
\jdanswer{
\begin{enumerate}
	\item
	6.010，4.20；
	\item
	①1.0；②参见下图；③1.66，充分利用已测得数据；④CD。
\end{enumerate}
}
%% memo:
\memoanswer{
（1）刻度尺的最小刻度为 $1\Umm$，所以读数时需要估读，根据图甲可知，金属杆的长度为 $60.10\Umm$（即 $6.010\Ucm$）；游标卡尺是不需要估读的，此卡尺的分度值是 $0.02\Umm$，所以金属杆的直径为 $4\Umm + 10 \times 0.02\Umm = 4.20\Umm$。

（2）①改装电流表扩大量程 $I = I_{g} + \frac{I_{g} r}{R}$，$R = \frac{I_{g} r}{I - I_{g}} = 1.0\UO$。②对照电路图，逐个回路连接即可，如图所示。

\onepicture[w=5.50cm]{figs/19_an.jpg}

③$r = \frac{\frac{U_{5} - U_{1}}{I_{5} - I_{1}} + \frac{U_{6} - U_{2}}{I_{6} - I_{2}} + \frac{U_{7} - U_{3}}{I_{7} - I_{3}} + \frac{U_{8} - U_{4}}{I_{8} - I_{4}}}{4} = 1.66\UO$，这样做的优点可以利用每一组数据。④电压表的读数等于路端电压，所以电压表内阻对测量结果没有影响，选项 A 错误；滑动变阻器用来改变路端电压，对测量结果没有影响，选项 B 错误；若 $R_{1}$ 的实际阻值偏小，则电流表的实际量程偏大，电流表的读数比实际值偏小，电池组的内阻测量值偏大，选项 C 正确；$R_{0}$ 的实际阻值偏大时，电池组的实际内阻比测量值偏小，选项 D 正确。
}

\item
%% number: 20
%% paperName: 2014年福建理综第 20 题 15 分
%% typeId: 6
%% type: 计算题
%% chapter: 电场
%% point: 电场叠加
%% method: 无
%% score: 15
%% degree: 650
%% duplicateId: 0
%% body:
如图，真空中 $xOy$ 平面直角坐标系上的 ABC 三点构成等边三角形，边长 $L = 2.0\Um$。若将电荷量均为 $q = +2.0\times10^{-6}\UC$ 的两点电荷分别固定在 A、B 点，已知静电力常量 $k = 9.0\times10^{9}\UNmqCq$。求：

\begin{enumerate}
	\item
	两点电荷间的库仑力大小；
	\item
	C 点的电场强度的大小和方向。
\end{enumerate}

\onepicture[w=5.50cm, align=right]{figs/20.jpg}
%% answer: 
\jdanswer{
\begin{enumerate}
	\item
	$9.0\times10^{-3}\UN$
	\item
	$7.8\times10^{3}\UNC$，方向沿 $y$ 轴正方向
\end{enumerate}
}
%% memo:
\memoanswer{
（1）根据库仑定律，A、B 两点电荷间的库仑力的大小为 $F = k \frac{q^{2}}{L^{2}}$，代入数据得 $F = 9.0 \times 10^{-3}\UN$。

（2）A、B 两点电荷在 C 点产生的场强大小相等，均为 $E_{1}=k\frac{q}{L^{2}}$；A、B 两点电荷形成的电场在 C 点的合场强大小为 $E_{\text{合}} = 2E_{1}\cos30^{\circ}$，代入数据得 $E_{\text{合}} = 7.8 \times 10^{3}\UNC$，方向沿 $y$ 轴正方向。
}

\item
%% number: 21
%% paperName: 2014年福建理综第 21 题 19 分
%% typeId: 6
%% type: 计算题
%% chapter: 曲线运动
%% point: 竖直平面圆周运动
%% method: 无
%% score: 19
%% degree: 450
%% duplicateId: 0
%% body:
图为某游乐场内水上滑梯轨道示意图，整个轨道在同一竖直平面内，表面粗糙的 AB 段与四分之一光滑圆弧轨道 BC 在 B 点水平相切。点 A 距水面的高度为 $H$，圆弧轨道 BC 的半径为 $R$，圆心 O 恰在水面。一质量为 $m$ 的游客（视为质点）可从轨道 AB 的任意位置滑下，不计空气阻力。

\begin{enumerate}
	\item
	若游客从 A 点由静止开始滑下，到 B 点时沿切线方向滑离轨道落在水面 D 点，$OD=2R$，求游客滑到 B 点时的速度 $v_{B}$ 大小及运动过程轨道摩擦力对其所做的功 $W_{f}$；
	\item
	若游客从 AB 段某处滑下，恰好停在 B 点，又因为受到微小扰动，继续沿圆弧轨道滑到 P 点后滑离轨道，求 P 点离水面的高度 $h$。（提示：在圆周运动过程中任一点，质点所受的向心力与其速率的关系为 $F_{\text{向}} = m\frac{v^{2}}{R}$）
\end{enumerate}

\onepicture[w=7.00cm, align=right]{figs/21.jpg}
%% answer: 
\jdanswer{
\begin{enumerate}
	\item
	$v_{B}=\sqrt{2gR}$，$W_{f}=2mgR-mgH$
	\item
	$h=\frac{2}{3}R$
\end{enumerate}
}
%% memo:
\memoanswer{
（1）游客从 B 点做平抛运动，有 $2R = v_{B}t$，$R = \frac{1}{2}gt^{2}$，联立解得 $v_{B} = \sqrt{2gR}$。从 A 到 B，根据动能定理，有 $mg(H - R) + W_{f} = \frac{1}{2}mv_{B}^{2}$，可得 $W_{f} = -(mgH - 2mgR)$。

（2）设 OP 与 OB 间夹角为 $\theta$，游客在 P 点时的速度为 $v_{P}$，受到的支持力为 $N$，从 B 到 P 由机械能守恒定律，有 $mg(R - R \cos \theta) = \frac{1}{2}mv_{P}^{2} - 0$。过 P 点根据向心力公式，有 $mg \cos \theta - N = \frac{mv_{P}^{2}}{R}$，$N = 0$，$\cos \theta = \frac{h}{R}$，解得 $h = \frac{2R}{3}$。
}

\item
%% number: 22
%% paperName: 2014年福建理综第 22 题 20 分
%% typeId: 6
%% type: 计算题
%% chapter: 磁场
%% point: 带电粒子在磁场中的运动
%% method: 无
%% score: 20
%% degree: 450
%% duplicateId: 0
%% body:
如图，某一新型发电装置的发电管是横截面为矩形的水平管道，管道的长为 $L$、宽度为 $d$、高为 $h$，上下两面是绝缘板，前后两侧面 M、N 是电阻可忽略的导体板，两导体板与开关 S 和定值电阻 $R$ 相连。整个管道置于磁感应强度大小为 $B$，方向沿 $z$ 轴正方向的匀强磁场中。管道内始终充满电阻率为 $\rho$ 的导电液体（有大量的正、负离子），且开关闭合前后，液体在管道进、出口两端压强差的作用下，均以恒定速率 $v_{0}$ 沿 $x$ 轴正向流动，液体所受的摩擦阻力不变。

\begin{enumerate}
	\item
	求开关闭合前，M、N 两板间的电势差大小 $U_{0}$；
	\item
	求开关闭合前后，管道两端压强差的变化 $\Delta p$；
	\item
	调整矩形管道的宽和高，但保持其它量和矩形管道的横截面 $S=dh$ 不变，求电阻 $R$ 可获得的最大功率 $P_{m}$ 及相应的宽高比 $d/h$ 的值。
\end{enumerate}

\onepicture[w=7.00cm, align=right]{figs/22.png}
%% answer: 
\jdanswer{
\begin{enumerate}
	\item
	$U_{0} = Bdv_{0}$
	\item
	$\Delta p = \frac{Ldv_{0}B^{2}}{LhR + d\rho}$
	\item
	$P_{m} = \frac{LSv_{0}^{2}B^{2}}{4\rho}$，$\frac{d}{h}=\frac{LR}{\rho}$
\end{enumerate}
}
%% memo:
\memoanswer{
（1）设带电粒子所带的电量为 $q$，当其所受的洛伦兹力与电场力平衡时，$U_{0}$ 保持恒定，有 $qv_{0}B = q\frac{U_{0}}{d}$，得 $U_{0} = Bdv_{0}$。

（2）设开关闭合前后，管道两端压强差分别为 $p_{1}$、$p_{2}$，液体所受摩擦阻力均为 $f$，开关闭合后管道内液体受到安培力为 $F_{\text{安}}$，有 $p_{1}hd = f$，$p_{2}hd = f + F_{\text{安}}$，$F_{\text{安}} = BId$；根据欧姆定律，有 $I = \frac{U_{0}}{R + r}$；两板间液体的电阻 $r = \rho \frac{d}{Lh}$；联立解得：$\Delta p = \frac{Ldv_{0}B^{2}}{LhR + d\rho}$。

（3）电阻 R 获得的功率为 $P = I^{2}R$，$P = \frac{Lv_{0}B}{\frac{LR}{d} + \frac{\rho}{h}} R$，当 $\frac{d}{h} = \frac{LR}{\rho}$ 时，电阻 R 获得最大功率 $P_{m} = \frac{LSv_{0}^{2}B^{2}}{4\rho}$。
}

\item
%% number: 29
%% paperName: 2014年福建理综第 29 题 5 分
%% typeId: 1
%% type: 单选题
%% chapter: 气体性质
%% point: 气体图象
%% method: 无
%% score: 5
%% degree: 650
%% duplicateId: 0
%% body:
【物理——选修 3-3】（本题共有两小题，每小题 6 分，共 12 分。每小题只有一个选项符合题意）

（1）如图，横坐标 $v$ 表示分子速率，纵坐标 $f(v)$ 表示各等间隔速率区间的分子数占总分子数的百分比。图中曲线能正确表示某一温度下气体分子麦克斯韦速率分布规律的是\tkanswer[1.2cm]{D}。（填选项前的字母）

\onepicture[w=5.50cm]{figs/29a.png}

\fourchoices
{曲线①}
{曲线②}
{曲线③}
{曲线④}

（2）图为一定质量理想气体的压强 $p$ 与体积 $V$ 关系图像，它由状态 A 经等容过程到状态 B，再经等压过程到状态 C。设 A、B、C 状态对应的温度分别为 $T_{A}$、$T_{B}$、$T_{C}$，则下列关系式中正确的是\xzanswer{C}。（填选项前的字母）

\onepicture[w=4.50cm]{\includesvg{figs/29b}}

\fourchoices[answer=C]
{$T_{A} < T_{B}$，$T_{B} < T_{C}$}
{$T_{A} > T_{B}$，$T_{B} = T_{C}$}
{$T_{A} > T_{B}$，$T_{B} < T_{C}$}
{$T_{A} = T_{B}$，$T_{B} > T_{C}$}
%% answer: C
%% memo:
\memoanswer{
（1）气体分子速率分布是有规律的，速率较大或较小的分子占少数，接近平均速率的分子占多数，分子速率不可能为 0，也不可能为无穷大，因此只有曲线④符合要求。D 项正确。

（2）一定质量的理想气体，由状态 A 经等容过程到状态 B 有 $\frac{p_{A}}{T_{A}} = \frac{p_{B}}{T_{B}}$，由于 $p_{A} > p_{B}$，所以 $T_{A} > T_{B}$；由状态 B 经等压过程到状态 C 有 $\frac{V_{B}}{T_{B}} = \frac{V_{C}}{T_{C}}$，由于 $V_{B} < V_{C}$，所以 $T_{B} < T_{C}$。故 C 项正确。
}

\item
%% number: 30
%% paperName: 2014年福建理综第 30 题 5 分
%% typeId: 1
%% type: 单选题
%% chapter: 运动和力
%% point: 动量守恒定律
%% method: 无
%% score: 5
%% degree: 650
%% duplicateId: 0
%% body:
【物理——选修 3-5】（本题共有 2 小题，每小题 6 分，共 12 分。每小题只有一个选项符合题意）

（1）如图，放射性元素镭衰变过程中释放出 $\alpha$、$\beta$、$\gamma$ 三种射线，分别进入匀强电场和匀强磁场中，下列说法正确的是\tkanswer[1.2cm]{C}。（填选项前的字母）

\onepicture[w=5.50cm]{figs/30a.jpg}

\fourchoices
{①表示 $\gamma$ 射线，③表示 $\alpha$ 射线}
{②表示 $\beta$ 射线，③表示 $\alpha$ 射线}
{④表示 $\alpha$ 射线，⑤表示 $\gamma$ 射线}
{⑤表示 $\beta$ 射线，⑥表示 $\alpha$ 射线}

（2）一枚火箭搭载着卫星以速率 $v_{0}$ 进入太空预定位置，由控制系统使箭体与卫星分离。已知前部分的卫星质量为 $m_{1}$，后部分的箭体质量为 $m_{2}$，分离后箭体以速率 $v_{2}$ 沿火箭原方向飞行，若忽略空气阻力及分离前后系统质量的变化，则分离后卫星的速率 $v_{1}$ 为\xzanswer{D}。（填选项前的字母）

\onepicture[w=7.00cm]{figs/30b.png}

\fourchoices[answer=D]
{$v_{0} - v_{2}$}
{$v_{0} + v_{2}$}
{$v_{0} - \frac{m_{2}}{m_{1}}v_{2}$}
{$v_{0} + \frac{m_{2}}{m_{1}}$（$v_{0}-v_{2}$）}
%% answer: D
%% memo:
\memoanswer{
（1）$\alpha$ 射线带正电，$\beta$ 射线带负电，$\gamma$ 射线不带电。在匀强电场中，$\alpha$ 射线与 $\beta$ 射线分别在电场力的作用下发生偏转，$\alpha$ 射线偏向负极板，$\beta$ 射线偏向正极板，$\gamma$ 射线不受电场力，不发生偏转；在磁场中，由左手定则可以判断 $\alpha$ 射线向左偏，$\beta$ 射线向右偏，$\gamma$ 射线不受洛伦兹力，不发生偏转。故 C 项正确。

（2）忽略空气阻力和分离前后系统质量的变化，卫星和箭体整体分离前后动量守恒，根据动量守恒定律则有 $(m_{1} + m_{2})v_{0} = m_{2}v_{2} + m_{1}v_{1}$，解得 $v_{1} = v_{0} + \frac{m_{2}}{m_{1}}(v_{0} - v_{2})$；D 项正确。
}

\end{enumerate}

\end{document}
